%% toc: P3 34 34.9
\chapter{Real Closed Fields as the Principal Worked Example}
\label{sec:p03-w03-p03-037-rcf}

The theory of real closed fields is the example against which the notions of
\cref{sec:p03-nip-ominimal} are usually tested.  It is unstable, because its
order is a definable linear order; it is o-minimal, because every
one-variable definable set is a finite union of points and intervals; and it
is NIP.  All three facts rest on one theorem: in the language of ordered
rings the theory eliminates quantifiers, so that the definable sets of a real
closed field are exactly its semialgebraic sets.  This chapter proves that
theorem from the axioms and derives the rest from it.

\emph{Purpose.}  To give one theory in which the definitions of o-minimality,
NIP and stability can be checked by hand, with every hypothesis visible: the
language, the axioms, the parameter conventions and the level (one model or
the whole theory) at which each property holds.

\emph{Inputs.}  The ordered-ring language, the ordered-field axioms, and the
intermediate-value scheme for polynomials.  The model-theoretic notions are
those of \cref{def:p03-w03-p03-037-qe-def,def:p03-nip,def:p03-o-minimal,def:p03-ominimal-cell,def:p03-stability}.

\emph{Conclusions.}  Quantifier elimination
(\cref{thm:p03-w03-p03-037-qe}); completeness
(\cref{thm:p03-w03-p03-037-complete}); definable sets are semialgebraic
(\cref{thm:p03-w03-p03-037-definable-semialgebraic}); every model is
o-minimal (\cref{thm:p03-w03-p03-037-o-minimal}); explicit monotonicity and
cell-decomposition instances (\cref{def:p03-w03-p03-037-monotone-example,def:p03-w03-p03-037-circle-cells}); instability
(\cref{thm:p03-w03-p03-037-unstable}); a uniform bound for one-variable
formulas (\cref{thm:p03-w03-p03-037-nip-one-variable}) and theory-level NIP
for formulas with object tuples of any length
(\cref{thm:p03-w03-p03-037-nip}); and two boundary results, that elimination
fails without the order symbol (\cref{thm:p03-w03-p03-037-ring-language}) and
that one added predicate destroys both o-minimality and NIP
(\cref{thm:p03-w03-p03-037-expansion}).  Every result is proved in
Annex~\ref{sec:annex-p03-w03-p03-037-rcf}; the general o-minimal monotonicity
and cell decomposition theorems are not reproved here, and the worked
instances below check their conclusions directly.

\section{Language, axioms and semialgebraic sets}
\label{sec:p03-w03-p03-037-axioms}

\begin{hbdefinition}{def:p03-w03-p03-037-rcf}{Real closed fields}
The \emph{ordered-ring language} is
\(\Lang_{\mathrm{or}}=\{0,1,+,-,\mathbin{\cdot},<\}\), with \(-\) binary
subtraction.  An \emph{ordered field} is an \(\Lang_{\mathrm{or}}\)-structure
that is a commutative field whose relation \(<\) is a strict linear order
satisfying \(x<y\rightarrow x+z<y+z\) and
\((0<x\wedge 0<y)\rightarrow 0<x\mathbin{\cdot}y\).
For \(n\geq 1\) let \(p_n(a_0,\dots,a_n,t)\) be the term
\(a_n t^n+\dots+a_1t+a_0\) (powers written as iterated products).  The
theory \(\mathsf{RCF}\) consists of the ordered-field axioms together with,
for every \(n\geq1\), the \emph{intermediate-value axiom}
\begin{multline*}
\forall a_0\dots\forall a_n\,\forall u\,\forall v\,
\bigl((u<v\wedge p_n(\bar a,u)\mathbin{\cdot}p_n(\bar a,v)<0)\\ \rightarrow\exists x\,(u<x\wedge x<v\wedge p_n(\bar a,x)=0)\bigr).
\end{multline*}
A model of \(\mathsf{RCF}\) is a \emph{real closed field}.  In such a field
\(K\), \(\operatorname{sgn}:K\to\{-1,0,1\}\) is the sign function and
\(\abs{x}\) is \(x\) or \(-x\), whichever is not negative.
\index{real closed field}
\end{hbdefinition}

The intermediate-value form is one of three equivalent characterizations of
real closed fields listed by \cite[p.~8]{Coste1999OMin}; the other two are
that every positive element is a square and every odd-degree polynomial has
a root, and that adjoining a square root of \(-1\) gives an algebraically
closed field.  Those equivalences are stated there without proof and are not
used below: every result of this chapter is derived from the axioms of
\cref{def:p03-w03-p03-037-rcf} as displayed.  The field \(\mathbb R\) of real
numbers is a real closed field by the intermediate-value theorem of real
analysis, and the field of real algebraic numbers is the smallest real closed
field \cite[p.~9]{Coste1999OMin}.  Every ordered field has
characteristic \(0\), since \(n\cdot1\) is a sum of \(n\) positive elements;
hence the integers embed in it as \(n\mapsto n\cdot 1\).

\begin{hbconvention}{def:p03-w03-p03-037-parametric}{Integer polynomials and parametric polynomials}
Every \(\Lang_{\mathrm{or}}\)-term \(t(\bar z)\) denotes, in every ordered
field, the function of an integer polynomial \(P_t\in\mathbb Z[\bar z]\), and
every integer polynomial is denoted by some term.  We therefore write
\(P=0\), \(P>0\) and \(P<0\) for the atomic formulas \(t=0\), \(0<t\) and
\(t<0\) with \(P_t=P\).  A \emph{parametric polynomial} in the parameter
variables \(\bar Y=(Y_1,\dots,Y_m)\) is a polynomial
\(P=\sum_{k\leq p}a_kX^k\) with coefficients \(a_k\in\mathbb Z[\bar Y]\)
together with its \emph{formal degree} \(p\geq0\); its \emph{formal leading
coefficient} \(a_p\) may be the zero polynomial.  For \(j\leq p\) the
\emph{truncation} \(P_{\leq j}\) is \(\sum_{k\leq j}a_kX^k\) with formal
degree \(j\), and the \emph{formal derivative} \(P'\) is
\(\sum_{1\leq k\leq p}k\,a_kX^{k-1}\) with formal degree \(p-1\) (for
\(p\geq1\)).  For a real closed field \(K\) and \(\bar y\in K^m\),
\(P(\bar y)\in K[X]\) is obtained by evaluating the coefficients at
\(\bar y\).
\end{hbconvention}

\begin{hbdefinition}{def:p03-w03-p03-037-semialgebraic}{Semialgebraic sets}
Let \(K\) be a real closed field.  A set \(S\subseteq K^n\) is
\emph{semialgebraic} when it is a finite Boolean combination of sets
\(\{\bar x\in K^n:P(\bar x)>0\}\) and \(\{\bar x\in K^n:P(\bar x)=0\}\) with
\(P\in K[X_1,\dots,X_n]\).  For \(K=\mathbb R\) this is the class of
\cite[p.~5]{Coste1999OMin}, and the same definition is used there over an
arbitrary real closed field \cite[p.~9]{Coste1999OMin}.
\index{semialgebraic set}
\end{hbdefinition}

\begin{hbdefinition}{def:p03-w03-p03-037-qe-def}{Quantifier elimination}
A theory \(T\) in a signature \(\Lang\) \emph{eliminates quantifiers} when
for every \(\Lang\)-formula \(\varphi(\bar x)\) there is a quantifier-free
\(\Lang\)-formula \(\psi(\bar x)\) with the same free variables such that
\(T\entails\forall\bar x\,(\varphi(\bar x)\leftrightarrow\psi(\bar x))\).
The equivalence is required in every model of \(T\), not in every
\(\Lang\)-structure, and the property depends on the signature: enlarging
or shrinking \(\Lang\) can create or destroy it.
\index{quantifier elimination}
\end{hbdefinition}

\begin{hbwarning}{def:p03-w03-p03-037-language}{The order symbol is a hypothesis}
Every result of this chapter is stated for \(\Lang_{\mathrm{or}}\).  In the
ring language \(\{0,1,+,-,\mathbin{\cdot}\}\) the same real closed fields
have the same parameter-definable sets, because \(x\leq y\) is defined by
\(\exists z\,(y-x=z\mathbin{\cdot}z)\); but quantifier elimination fails there
(\cref{thm:p03-w03-p03-037-ring-language}), and the definition of
o-minimality in \cref{def:p03-o-minimal} applies to an expansion of a dense
linear order, which a pure field is not until its order is named.
\end{hbwarning}

\section{Elementary analysis in a real closed field}
\label{sec:p03-w03-p03-037-analysis}

The elimination argument uses only the following facts about one-variable
polynomials, each proved from the intermediate-value axioms.

\begin{hblemma}{thm:p03-w03-p03-037-poly-basics}{proved=annex:p03-w03-p03-037-poly-basics}{One-variable polynomial facts}
Let \(K\) be a real closed field and \(P=\sum_{k\leq p}a_kX^k\in K[X]\) with
\(a_p\neq0\).
\begin{enumerate}
\item \(P\) has at most \(p\) roots in \(K\).
\item If \(C\subseteq K\) is convex and contains no root of \(P\), then
  \(\operatorname{sgn}P\) is constant on \(C\).
\item Put
  \[M_P=1+\sum_{k<p}\abs{a_k}/\abs{a_p}.\]
  If \(\abs{x}>M_P\) then
  \(\operatorname{sgn}P(x)=\operatorname{sgn}(a_p)\operatorname{sgn}(x)^p\).
\item (Rolle) If \(a<b\) and \(P(a)=P(b)=0\), then \(P'\) has a root in
  \((a,b)\).
\item (Mean value) If \(a<b\), there is \(c\in(a,b)\) with
  \(P(b)-P(a)=(b-a)P'(c)\).
\item (Monotonicity) Let \(J\) be an open interval with endpoints in
  \(K\cup\{\pm\infty\}\), let \(p\geq1\), and suppose \(P'\) has no root in
  \(J\).  Then \(P'\) has a constant sign \(\varepsilon\neq0\) on \(J\), and
  \(\varepsilon\,(P(x')-P(x))>0\) whenever \(x<x'\) lie in \(J\) or are
  endpoints of \(J\) belonging to \(K\).
\item Every polynomial function \(K\to K\) is continuous for the order
  topology, and every \(x\geq0\) has exactly one square root
  \(\sqrt{x}\geq0\); moreover \(\abs{\sqrt u-\sqrt v}\leq\sqrt{\abs{u-v}}\)
  for \(u,v\geq0\), so \(\sqrt{\cdot}\) is continuous on \([0,\infty)\).
\end{enumerate}
\end{hblemma}

\begin{hblemma}{thm:p03-w03-p03-037-interval}{proved=annex:p03-w03-p03-037-interval}{Behaviour of a polynomial between roots of its derivative}
Let \(K\) be a real closed field, \(P\in K[X]\) of degree \(p\geq1\) with
leading coefficient \(a_p\), and \(J=(\alpha,\beta)\) with
\(\alpha<\beta\) in \(K\cup\{\pm\infty\}\), such that \(P'\) has no root in
\(J\); let \(\varepsilon\) be the sign of \(P'\) on \(J\).  Put
\(s_\alpha=\operatorname{sgn}P(\alpha)\) if \(\alpha\in K\) and
\(s_\alpha=(-1)^p\operatorname{sgn}a_p\) if \(\alpha=-\infty\); put
\(s_\beta=\operatorname{sgn}P(\beta)\) if \(\beta\in K\) and
\(s_\beta=\operatorname{sgn}a_p\) if \(\beta=+\infty\).  Then exactly one
of the following holds.
\begin{enumerate}
\item \(s_\alpha=0\), and \(P\) has sign \(\varepsilon\) at every point of \(J\).
\item \(s_\beta=0\), and \(P\) has sign \(-\varepsilon\) at every point of \(J\).
\item \(s_\alpha=s_\beta\neq0\), and \(P\) has sign \(s_\alpha\) at every point
  of \(J\).
\item \(s_\alpha=-s_\beta\neq0\), and \(P\) has exactly one root \(c\) in
  \(J\), with sign \(s_\alpha\) on \((\alpha,c)\) and \(s_\beta\) on
  \((c,\beta)\).
\end{enumerate}
\end{hblemma}

\section{Sign tables and quantifier elimination}
\label{sec:p03-w03-p03-037-qe}

\begin{hbdefinition}{def:p03-w03-p03-037-sign-table}{Sign table of a list of polynomials}
Let \(K\) be a real closed field and \(\mathcal P=(P_1,\dots,P_s)\) a list in
\(K[X]\).  Let \(z_1<\dots<z_N\) be the roots in \(K\) of those \(P_i\) that
are not the zero polynomial.  The \emph{cells} of \(\mathcal P\) are
\[
I_0=(-\infty,z_1),\ \{z_1\},\ I_1=(z_1,z_2),\ \dots,\ \{z_N\},\
I_N=(z_N,+\infty),
\]
with \(I_0=K\) when \(N=0\).  They partition \(K\), each is nonempty, and
each \(P_i\) has constant sign on each cell
(\cref{thm:p03-w03-p03-037-poly-basics}); \(\sigma(C)\in\{-1,0,1\}^s\) is the
vector of these signs on the cell \(C\).  The \emph{sign table} is the word
\[
\operatorname{SIGN}(\mathcal P)=
\bigl(\sigma(I_0),\sigma(\{z_1\}),\sigma(I_1),\dots,
\sigma(\{z_N\}),\sigma(I_N)\bigr)
\]
of length \(2N+1\) over the alphabet \(\{-1,0,1\}^s\).  For a list of
parametric polynomials and \(\bar y\in K^m\), \(\mathcal P(\bar y)\) is the
list \((P_1(\bar y),\dots,P_s(\bar y))\).  If the formal degrees sum to
\(D\), then \(N\leq D\), so \(\operatorname{SIGN}(\mathcal P(\bar y))\)
ranges over the finite set \(\mathbb T_{s,D}\) of words of length at most
\(2D+1\).
\end{hbdefinition}

The sign table is a finite combinatorial object that does not mention \(K\).
The elimination theorem says that which word it is depends on the parameters
through a quantifier-free condition, uniformly over all real closed fields.

\begin{hblemma}{thm:p03-w03-p03-037-sublist}{proved=annex:p03-w03-p03-037-sublist}{Sign tables of sublists}
For every injective map \(\iota\) from \(\{1,\dots,r\}\) to \(\{1,\dots,s\}\) there is a
function \(\rho_\iota\) on words, depending only on \(\iota\), such that for
every real closed field and every list \(\mathcal P=(P_1,\dots,P_s)\),
\(\operatorname{SIGN}(P_{\iota(1)},\dots,P_{\iota(r)})=
\rho_\iota(\operatorname{SIGN}(\mathcal P))\).
\end{hblemma}

\begin{hbmethod}{def:p03-w03-p03-037-elimination-step}{One elimination step}
Let \(\mathcal P=(P_1,\dots,P_{s-1},P)\) be a list of parametric polynomials
in which \(P=\sum_{k\leq p}a_kX^k\) has the largest formal degree
\(p\geq1\), and let \(\mathcal P_0=(P_1,\dots,P_{s-1})\).
\begin{enumerate}
\item \emph{Degree drop.}  Where \(a_p(\bar y)=0\), replace \(P\) by its
  truncation \(P_{\leq p-1}\); the list \(\mathcal P^\flat\) so obtained has
  the same sign table at \(\bar y\).
\item \emph{Auxiliary list.}  Let \(\mathcal R=(P_1,\dots,P_{s-1},P')\).  The
  list \(\mathcal Q\) consists of \(\mathcal P_0\), \(P'\), the constant
  \(a_p\), every coefficient \(b_k\) of every member \(Q\) of \(\mathcal R\)
  (as a polynomial of formal degree \(0\)), and, for every member \(Q\) of
  \(\mathcal R\) of formal degree \(q\) and every \(1\leq j\leq q\), the
  pseudo-remainder \(\operatorname{prem}(P,Q_{\leq j})\) of
  \cref{def:p03-w03-p03-037-prem}.
\item \emph{Reading off.}  Where \(a_p(\bar y)\neq0\), the sign table of
  \(\mathcal P(\bar y)\) is computed from that of \(\mathcal Q(\bar y)\):
  the roots of \(\mathcal R\) are located, the sign of \(P\) at each of them
  is read from a pseudo-remainder, the signs at \(\pm\infty\) from \(a_p\),
  and each interval between consecutive roots of \(\mathcal R\) is resolved by
  \cref{thm:p03-w03-p03-037-interval}.
\end{enumerate}
Side conditions: the pseudo-remainder used at a root \(w\) of
\(Q(\bar y)\) is the one for the truncation \(Q_{\leq j}\) at the true degree
\(j\) of \(Q(\bar y)\), which the signs of the coefficients \(b_k\) reveal;
\(j\geq1\) because a nonzero constant has no root.  Every new member of
\(\mathcal Q\) has formal degree below \(p\), so the step decreases the pair
(largest formal degree, number of members of that degree) lexicographically.
The reduction of a parametric system to fixed degrees by truncation is the
one used by \cite[p.~11]{Coste2002SAG}.
\end{hbmethod}

\begin{hbdefinition}{def:p03-w03-p03-037-prem}{Pseudo-remainder}
Let \(P\) and \(Q\) be parametric polynomials of formal degrees \(p\geq q\geq1\),
and let \(b_q\) be the formal leading coefficient of \(Q\).  The
\emph{pseudo-remainder} \(\operatorname{prem}(P,Q)\) is the parametric
polynomial \(R\) of formal degree \(q-1\) produced by the recursion of
\cref{thm:p03-w03-p03-037-prem} below, which satisfies
\(b_q^{\,p-q+1}P=A\,Q+R\) in \(\mathbb Z[\bar Y][X]\) for some \(A\).
\end{hbdefinition}

\begin{hblemma}{thm:p03-w03-p03-037-prem}{proved=annex:p03-w03-p03-037-prem}{Pseudo-division}
Let \(Q=\sum_{k\leq q}b_kX^k\) have formal degree \(q\geq1\).  For every
\(P\) of formal degree \(p\geq q-1\) there are \(A,R\in\mathbb Z[\bar Y][X]\),
\(R\) of formal degree \(q-1\), with \(b_q^{\,p-q+1}P=A\,Q+R\).  Consequently,
if \(K\) is a real closed field, \(\bar y\in K^m\), \(b_q(\bar y)\neq0\) and
\(Q(\bar y)(w)=0\), then
\(\operatorname{sgn}P(\bar y)(w)=
\operatorname{sgn}(b_q(\bar y))^{p-q+1}\operatorname{sgn}R(\bar y)(w)\).
\end{hblemma}

\begin{hbtheorem}{thm:p03-w03-p03-037-uniform-sign}{proved=annex:p03-w03-p03-037-uniform-sign}{Uniform sign tables}
Let \(\mathcal P\) be a finite list of parametric polynomials in the
parameter variables \(\bar Y\).  For every word \(t\) there is a
quantifier-free \(\Lang_{\mathrm{or}}\)-formula \(\theta_{\mathcal P,t}(\bar Y)\)
such that for every real closed field \(K\) and every \(\bar y\in K^m\),
\[
\operatorname{SIGN}(\mathcal P(\bar y))=t
\iff K\entails\theta_{\mathcal P,t}(\bar y).
\]
\end{hbtheorem}

\begin{hbtheorem}{thm:p03-w03-p03-037-qe}{proved=annex:p03-w03-p03-037-qe}{Quantifier elimination for real closed fields}
\(\mathsf{RCF}\) eliminates quantifiers in \(\Lang_{\mathrm{or}}\) in the
sense of \cref{def:p03-w03-p03-037-qe-def}: for every
\(\Lang_{\mathrm{or}}\)-formula \(\varphi(\bar x)\) there is a quantifier-free
\(\psi(\bar x)\) with
\(\mathsf{RCF}\entails\forall\bar x\,(\varphi\leftrightarrow\psi)\).
\end{hbtheorem}

This is the Tarski--Seidenberg elimination; in the form of
\cite[p.~11, Theorem~1.9]{Coste2002SAG} it says that the existence of a
real solution of a system of polynomial equations and inequalities with
parameters is equivalent to a finite disjunction of systems in the
parameters alone.  The proof in the annex is the sign-table argument above
and applies to every real closed field at once.

\begin{hbexample}{def:p03-w03-p03-037-quadratic}{Eliminating the variable of a quadratic}
In every real closed field,
\begin{multline*}
\exists x\,(ax^2+bx+c=0)\ \longleftrightarrow\
(a\neq0\wedge b^2-4ac\geq0)\\ \vee(a=0\wedge b\neq0)\vee(a=0\wedge b=0\wedge c=0),
\end{multline*}
the instance displayed by \cite[p.~11]{Coste2002SAG}.  For \(a\neq0\),
\(ax^2+bx+c=a\bigl((x+b/2a)^2-(b^2-4ac)/4a^2\bigr)\), so a root exists exactly
when \((b^2-4ac)/4a^2\) is a square, that is, when \(b^2-4ac\geq0\): squares
are nonnegative in an ordered field, and nonnegative elements have square
roots by \cref{thm:p03-w03-p03-037-poly-basics}.  For \(a=0\) the equation
\(bx+c=0\) is solvable exactly when \(b\neq0\) or \(c=0\).  The case split
on whether \(a=0\) is the degree drop of
\cref{def:p03-w03-p03-037-elimination-step}.
\end{hbexample}

\begin{hbcorollary}{thm:p03-w03-p03-037-complete}{proved=annex:p03-w03-p03-037-complete}{Completeness}
\(\mathsf{RCF}\) is complete.  In particular every real closed field is
elementarily equivalent to \(\mathbb R\).
\end{hbcorollary}

\section{Definable sets and o-minimality}
\label{sec:p03-w03-p03-037-definable}

\begin{hbtheorem}{thm:p03-w03-p03-037-definable-semialgebraic}{proved=annex:p03-w03-p03-037-definable-semialgebraic}{Definable sets are semialgebraic}
Let \(K\) be a real closed field and \(S\subseteq K^n\).  Then \(S\) is
definable in \((K;0,1,+,-,\mathbin{\cdot},<)\) with parameters from \(K\) if
and only if \(S\) is semialgebraic.
\end{hbtheorem}

\begin{hbtheorem}{thm:p03-w03-p03-037-o-minimal}{proved=annex:p03-w03-p03-037-o-minimal}{Every real closed field is o-minimal}
Every real closed field, as an \(\Lang_{\mathrm{or}}\)-structure, is
o-minimal in the sense of \cref{def:p03-o-minimal}.  Since \(\mathsf{RCF}\)
is complete, \(\mathsf{RCF}\) is an o-minimal theory.  Moreover, if a
one-variable formula \(\varphi(x;\bar y)\) is equivalent modulo
\(\mathsf{RCF}\) to a Boolean combination of sign conditions on
\(P_1,\dots,P_s\in\mathbb Z[x,\bar y]\) whose \(x\)-degrees sum to \(D\), then
in every real closed field and for every \(\bar b\) the set
\(\varphi(K;\bar b)\) is a union of at most \(D+1\) pairwise disjoint convex
sets.
\end{hbtheorem}

The o-minimality here is a property of the theory, not only of the field
\(\mathbb R\): the bound \(D+1\) depends on the formula and not on the model
or the parameters.  The general monotonicity and cell decomposition theorems,
\cref{thm:p03-ominimal-monotonicity,thm:p03-ominimal-cell-decomposition},
therefore apply to every real closed field.  The two examples below check
their conclusions directly for specific semialgebraic data, using only
\cref{thm:p03-w03-p03-037-poly-basics}.

\begin{hbexample}{def:p03-w03-p03-037-monotone-example}{Monotonicity for semialgebraic functions}
Let \(K\) be a real closed field.
\begin{enumerate}
\item Let \(f(x)=x^3-3x\) on \((-2,2)\).  Then \(f'(x)=3(x-1)(x+1)\) has the
  roots \(-1\) and \(1\), is positive on \((-2,-1)\) and \((1,2)\), and is
  negative on \((-1,1)\).  By the monotonicity clause of
  \cref{thm:p03-w03-p03-037-poly-basics}, the partition
  \((-2,-1),\{-1\},(-1,1),\{1\},(1,2)\) makes \(f\) strictly increasing,
  strictly decreasing and strictly increasing on the three intervals, and
  \(f\) is continuous everywhere.  No partition can omit the point \(-1\):
  \(f(x)-f(-1)=(x+1)^2(x-2)<0\) for every \(x\neq-1\) in \((-2,2)\), so
  \(f\) is not monotone on any interval having \(-1\) as an interior point.
\item Let \(g:(-1,1)\to K\) be \(g(x)=1/x\) for \(x\neq0\) and \(g(0)=0\).
  Its graph \(\{(x,y):-1<x<1\wedge((x\neq0\wedge xy=1)\vee(x=0\wedge y=0))\}\)
  is semialgebraic.  On \((-1,0)\) and on \((0,1)\), \(g\) is continuous and
  strictly decreasing, because \(1/x-1/x'=(x'-x)/(xx')>0\) when \(x<x'\) have
  the same sign, and \(\abs{1/x-1/c}\leq2\abs{x-c}/c^2\) when
  \(\abs{x-c}<\abs c/2\).  The partition \((-1,0),\{0\},(0,1)\) is the one
  the monotonicity theorem promises; \(g\) is not continuous at \(0\), which
  is why the theorem partitions into points and intervals rather than
  intervals alone.
\end{enumerate}
\end{hbexample}

\begin{hbexample}{def:p03-w03-p03-037-circle-cells}{A cell decomposition adapted to the circle and the disc}
Let \(K\) be a real closed field,
\(S=\{(x,y)\in K^2:x^2+y^2=1\}\) and \(B=\{(x,y)\in K^2:x^2+y^2<1\}\).  On
\((-1,1)\) let \(\xi(x)=\sqrt{1-x^2}\), which is definable, positive and
continuous by \cref{thm:p03-w03-p03-037-poly-basics}.  The thirteen cells of
\cref{def:p03-ominimal-cell}
\begin{itemize}
\item over \((-\infty,-1)\) and over \((1,+\infty)\): the full bands
  \((-\infty,+\infty)\);
\item over \(\{-1\}\) and over \(\{1\}\): the graph of the constant \(0\) and
  the bands \((-\infty,0)\) and \((0,+\infty)\);
\item over \((-1,1)\): the graphs of \(-\xi\) and \(\xi\) and the bands
  \((-\infty,-\xi)\), \((-\xi,\xi)\) and \((\xi,+\infty)\)
\end{itemize}
partition \(K^2\), and
\(S=\operatorname{graph}(-\xi)\cup\operatorname{graph}(\xi)\cup
\{(-1,0),(1,0)\}\) while \(B\) is the single band \((-\xi,\xi)\) over
\((-1,1)\).  Indeed, for \(\abs x<1\), \(x^2+y^2=1\) iff
\((y-\xi(x))(y+\xi(x))=0\), and \(x^2+y^2<1\) iff \(\abs y<\xi(x)\); for
\(x=\pm1\), \(x^2+y^2=1\) iff \(y=0\), and \(x^2+y^2<1\) never holds; for
\(\abs x>1\), \(x^2+y^2>1\).  This is the conclusion of
\cref{thm:p03-ominimal-cell-decomposition} for the family \(\{S,B\}\),
obtained by hand.
\end{hbexample}

\section{Instability, NIP, and the limits of the example}
\label{sec:p03-w03-p03-037-nip}

\begin{hbtheorem}{thm:p03-w03-p03-037-unstable}{proved=annex:p03-w03-p03-037-unstable}{Real closed fields are unstable}
The formula \(x<y\) has the order property in \(\mathsf{RCF}\)
(\cref{def:p03-stability}).  Hence \(\mathsf{RCF}\) is not stable.
\end{hbtheorem}

\begin{hbtheorem}{thm:p03-w03-p03-037-nip-one-variable}{proved=annex:p03-w03-p03-037-nip-one-variable}{One-variable formulas are NIP with an explicit bound}
Let \(\varphi(x;\bar y)\) be an \(\Lang_{\mathrm{or}}\)-formula with a single
object variable \(x\), and let \(D\) be as in
\cref{thm:p03-w03-p03-037-o-minimal} for a quantifier-free equivalent of
\(\varphi\).  In no real closed field does the family
\(\{\varphi(K;\bar b):\bar b\in K^{\abs{\bar y}}\}\) shatter a set of
\(2D+3\) elements.  In particular \(\varphi\) is NIP in \(\mathsf{RCF}\) in
the sense of \cref{def:p03-nip}.
\end{hbtheorem}

The one-variable bound does not by itself settle formulas \(\varphi(\bar x;\bar y)\)
whose object tuple \(\bar x\) has several coordinates: a set of pairs need not
be linearly ordered, so the convexity argument has nothing to act on.  For
those, elimination gives a different handle.  After elimination every
instance \(\varphi(K;\bar b)\) is a fixed Boolean combination of sets
\(\{\bar a:f(\bar a)>0\}\) where \(f\) ranges over a finite-dimensional space
of functions of \(\bar a\), and two combinatorial lemmas, valid over every
ordered field, bound how many subsets of a finite set such families can cut
out.

\begin{hbdefinition}{def:p03-w03-p03-037-trace}{Traces and shatter counts}
Let \(\mathcal F\) be a family of subsets of a set \(X\) and \(A\subseteq X\)
finite.  The \emph{trace} of \(\mathcal F\) on \(A\) is
\(\mathcal F\cap A=\{S\cap A:S\in\mathcal F\}\), and \(\mathcal F\)
\emph{shatters} \(A\) when \(\mathcal F\cap A\) is the set of all subsets of
\(A\).  For \(\varphi(\bar x;\bar y)\) and a model \(K\), the family
\(\{\varphi(K;\bar b):\bar b\in K^{\abs{\bar y}}\}\) shatters \(A\) in this
sense exactly when it does in the sense of \cref{def:p03-nip}.
\end{hbdefinition}

\begin{hblemma}{thm:p03-w03-p03-037-linear-shatter}{proved=annex:p03-w03-p03-037-linear-shatter}{Positivity sets in a finite-dimensional space}
Let \(F\) be an ordered field, \(X\) a set, and \(V\) a vector space over
\(F\) of functions \(X\to F\) (with pointwise operations) spanned by \(N\)
functions.  Then the family \(\{\{x\in X:f(x)>0\}:f\in V\}\) shatters no
subset of \(X\) with \(N+1\) elements.
\end{hblemma}

\begin{hblemma}{thm:p03-w03-p03-037-sauer-shelah}{proved=annex:p03-w03-p03-037-sauer-shelah}{Counting traces of a family of bounded shattering}
Let \(d\geq0\) and let \(\mathcal F\) be a family of subsets of \(X\) that
shatters no subset of \(X\) with \(d+1\) elements.  Then for every finite
\(A\subseteq X\) with \(n\) elements,
\[
\abs{\mathcal F\cap A}\leq\sum_{i=0}^{d}\binom ni\leq(n+1)^d .
\]
\end{hblemma}

\begin{hbtheorem}{thm:p03-w03-p03-037-nip}{proved=annex:p03-w03-p03-037-nip}{Real closed fields are NIP}
Let \(\varphi(\bar x;\bar y)\) be an \(\Lang_{\mathrm{or}}\)-formula, with
tuples \(\bar x\) and \(\bar y\) of any length.  Let \(\psi(\bar x;\bar y)\)
be a quantifier-free formula with
\(\mathsf{RCF}\entails\forall\bar x\,\forall\bar y\,(\varphi\leftrightarrow\psi)\)
(\cref{thm:p03-w03-p03-037-qe}), whose atomic formulas are, after the
normalization of the first step of that proof, sign conditions on integer
polynomials \(P_1,\dots,P_k\in\mathbb Z[\bar x,\bar y]\).  Write each
\(P_j=\sum_{\alpha\in I_j}q_{j,\alpha}(\bar y)\,\bar x^{\alpha}\) with
\(I_j\) a finite set of exponent vectors and
\(q_{j,\alpha}\in\mathbb Z[\bar y]\), let \(N_j=\abs{I_j}\) and
\(d=2\sum_{j\leq k}N_j\).  Then in no real closed field does the family
\(\{\varphi(K;\bar b):\bar b\in K^{\abs{\bar y}}\}\) shatter a set with at
least \(n_d=2^{d+1}(d+1)!\) elements.  Hence every
\(\Lang_{\mathrm{or}}\)-formula is NIP in \(\mathsf{RCF}\), and
\(\mathsf{RCF}\) is an NIP theory in the sense of \cref{def:p03-nip}.
\end{hbtheorem}

The bound \(n_d\) is crude; its role is to be independent of the model \(K\)
and of the parameters \(\bar b\), which is what NIP at the level of the theory
requires.  The proof uses quantifier elimination in \(\Lang_{\mathrm{or}}\) in
an essential way: it is the linear structure of the sets
\(\{\bar a:P(\bar a,\bar b)>0\}\) in the coefficients that is counted, and no
such structure is available for a formula that still contains quantifiers.

A second route reaches the same theorem from o-minimality alone.  Simon
proves that a partitioned formula is NIP exactly when the formula with object
and parameter variables exchanged is NIP
\cite[p.~14, Lemma~2.5]{Simon2015GuideManuscript}, and that a theory is NIP
as soon as every formula with a single parameter variable is NIP
\cite[p.~16, Proposition~2.11]{Simon2015GuideManuscript}; with the
one-variable bound of \cref{thm:p03-w03-p03-037-nip-one-variable} these give
NIP for \(\mathsf{RCF}\), and in the same way for every o-minimal theory,
which is how Simon lists the o-minimal case
\cite[p.~17, Example~2.12]{Simon2015GuideManuscript}.  That route needs
indiscernible sequences, Ramsey's theorem and compactness, and is not the one
proved here; the counting proof above needs only elimination and finite
combinatorics, but applies only where elimination to polynomial sign
conditions is available.

\begin{hbtheorem}{thm:p03-w03-p03-037-ring-language}{proved=annex:p03-w03-p03-037-ring-language}{Elimination needs the order symbol}
Let \(K\) be a real closed field.  The theory of the ring reduct
\((K;0,1,+,-,\mathbin{\cdot})\) does not eliminate quantifiers in the
language \(\{0,1,+,-,\mathbin{\cdot}\}\): the formula
\(\exists z\,(x=z\mathbin{\cdot}z)\) is not equivalent in \(K\) to any
quantifier-free formula of that language.
\end{hbtheorem}

\begin{hbtheorem}{thm:p03-w03-p03-037-expansion}{proved=annex:p03-w03-p03-037-expansion}{An expansion that is neither o-minimal nor NIP}
Let \(\mathcal R_{\mathbb Z}=(\mathbb R;0,1,+,-,\mathbin{\cdot},<,Z)\) where
\(Z\) is a unary predicate interpreted by the set of integers.  Then
\(\mathcal R_{\mathbb Z}\) is not o-minimal, and in its complete theory the
formula
\(\delta(x;y):=Z(x)\wedge Z(y)\wedge\exists w\,(Z(w)\wedge y=x\mathbin{\cdot}w)\)
has the independence property.
\end{hbtheorem}

\begin{hbwarning}{def:p03-w03-p03-037-expansions}{What the example does not show}
\Cref{thm:p03-w03-p03-037-o-minimal,thm:p03-w03-p03-037-nip} concern
\(\mathsf{RCF}\) in \(\Lang_{\mathrm{or}}\) and nothing larger.
\Cref{thm:p03-w03-p03-037-expansion} shows that adding one unary predicate to
the real field destroys both o-minimality and NIP, so neither property passes
to arbitrary expansions.  Particular expansions are o-minimal by separate,
case-specific theorems; Coste cites Wilkie's theorem that the real field with
the exponential function defines an o-minimal structure
\cite[p.~6]{Coste1999OMin}.  Conversely an ordered field admitting an
o-minimal expansion is already real closed \cite[p.~11,
Exercise~1.9]{Coste1999OMin}; that statement is recorded here only as a
source pointer and is not used.  Instability shows that o-minimality and NIP
do not imply stability; \(\mathsf{RCF}\) and \(\mathsf{DLO}\) are the
standing unstable NIP examples of the overview table in
\cref{sec:p03-nip-ominimal}.  The converse failure, an ordered theory that is
NIP but not o-minimal, is \cref{thm:p03-w03-p03-035-dense-predicate}, and the
placement of \(\mathsf{RCF}\) among the simple, stable and NIP classes is
\cref{thm:p03-w03-p03-035-rcf}.
\end{hbwarning}

\section{Relationships and sources}
\label{sec:p03-w03-p03-037-sources}

The chapter supplies the real-closed-field entries that the overview in
\cref{sec:p03-nip-ominimal} and the catalogue in
\cref{tab:p03-qe-catalog} record without proof: elimination in
\(\Lang_{\mathrm{or}}\), completeness, semialgebraic definability,
o-minimality of every model, instability and NIP.  The general o-minimal
monotonicity and cell decomposition theorems remain as stated at
\cref{thm:p03-ominimal-monotonicity,thm:p03-ominimal-cell-decomposition};
the examples here are instances checked from the definitions.  Sources read
for this chapter: \cite[pp.~5--11]{Coste1999OMin} for the definitions of
semialgebraic sets and the characterizations of real closed fields;
\cite[pp.~11--12]{Coste2002SAG} for the statement of Tarski--Seidenberg
elimination and the reduction to fixed degrees; and
\cite[pp.~13--17]{Simon2015GuideManuscript} for NIP, its symmetry and its
reduction to one variable, which this chapter records as an alternative
route rather than as a dependency.  Both Coste texts are the author's lecture notes
(dated November 1999 and October 2002) and Simon's text is his author-typeset
manuscript of 16 October 2015, not the printed book.
